%% -----------------------------------------------------------------
%% 北京邮电大学研究生学位论文 - 攻读学位期间取得的创新成果目录
%% 格式要求（2025年1月15日规范）：
%% 1. 标题为“攻读学位期间取得的创新成果目录”（注：最新规范不带双引号）
%% 2. 标题黑体三号字，居中对齐，首行无缩进，段前0行，段后2行，单倍行距
%% 3. 内容用宋体小四号，固定值20磅
%% 4. 按取得成果的时间顺序，列出本人在攻读学位期间取得的创新成果清单
%% -----------------------------------------------------------------

\begin{publication}

\section*{一、学术期刊与会议论文}
\begin{enumerate}[leftmargin=2.5em, itemsep=4pt, label={[\arabic*]}]
  \item \textbf{张三}, 李四, 王五. 基于自适应松弛机制的高效协同计算方法研究[J]. 软件学报, 2024, 35(8): 3610--3625. (CCF-A类中文核心期刊)
  \item \textbf{San Zhang}, Si Li, Wu Wang. Adaptive optimization framework for multi-tier edge networks[C]// Proceedings of the IEEE International Conference on Computer Communications (INFOCOM). IEEE, 2024: 1120--1129. (CCF-A类国际顶级会议)
  \item Si Li, \textbf{San Zhang}, Wu Wang. Robust distributed learning over fading channels[J]. IEEE Transactions on Mobile Computing, 2023, 22(11): 6580--6594. (CCF-A类国际期刊)
\end{enumerate}

\section*{二、申请与授权专利}
\begin{enumerate}[leftmargin=2.5em, itemsep=4pt, label={[\arabic*]}]
  \item 李四, \textbf{张三}, 王五. 一种面向边缘计算的动态自适应资源调度方法及系统: CN202310888888.8[P]. 2024-03-15. (实质审查阶段)
  \item \textbf{张三}, 李四. 一种高可靠异构网络数据传输优化方法: CN202210666666.6[P]. 2023-11-20. (已授权)
\end{enumerate}

\section*{三、参与科研项目}
\begin{enumerate}[leftmargin=2.5em, itemsep=4pt, label={[\arabic*]}]
  \item 国家自然科学基金重点项目：“面向下一代移动通信的高性能协同计算理论与关键技术”（项目编号：62231000），2022.01--2026.12，主要学生骨干参与。
  \item 国家重点研发计划课题：“异构网络自适应调度平台研制与应用示范”（项目编号：2023YFB2900000），2023.01--2025.12，主要参与人员。
\end{enumerate}

\end{publication}
